\documentclass[12pt]{article}
\usepackage[utf8]{inputenc}
\usepackage[spanish]{babel}
\usepackage{amsmath, amssymb}
\usepackage{booktabs}
\usepackage{geometry}
\usepackage{tikz}
\usepackage{pgfplots}
\pgfplotsset{compat=1.18}
\usetikzlibrary{fillbetween}
\usepackage{float}
\geometry{margin=2cm}
\usepackage{hyperref}

\usepackage{subcaption}

\begin{document}

\vspace{5cm}



\begin{center}
\begin{large}
\textbf{\\Guía 5: Econometría ICS-294}
\end{large}\\
\end{center}

\href{https://colab.research.google.com/drive/1PAX_qvkKfoZjpzktfvsFLToTcGMyIl69?usp=sharing}{LINK PAUTA en R}

\begin{enumerate}
    \item En una regresión lineal simple, el coeficiente asociado a la variable
        \(X_1\) fue estimado en:
        \[
          \hat\beta_1 = 2.7,
          \qquad s.e.(\hat\beta_1) = 0.9,
          \qquad n = 250.
        \]
        Considere un nivel de significancia del 5\%.

        \begin{enumerate}
        \item Plantee las hipótesis nula y alternativa para evaluar si
                \( \beta_1 \) es significativamente distinto de cero.
          \item Calcule el estadístico \( t \) de prueba, compare su valor
                con el valor crítico correspondiente y determine si se
                rechaza la hipótesis nula.
          \item Construya un intervalo de confianza al 95\% para
                \( \beta_1 \) e interprete su significado.
        \end{enumerate}

 {\color{red}       \textbf{Solución.}
        \begin{enumerate}
          \item Hipótesis a testear:
                \[
                  H_0: \beta_1 = 0
                  \qquad\text{versus}\qquad
                  H_1: \beta_1 \neq 0.
                \]
          \item El estadístico \( t \) es:
                \[
                  t
                  = \frac{\hat\beta_1 - 0}{s.e.(\hat\beta_1)}
                  = \frac{2.7}{0.9}
                  = 3.0.
                \]
                Con 248 grados de libertad (suponiendo pocos regresores)
                y \(n = 250\), podemos aproximar el valor crítico bilateral
                usando la normal estándar:
                \[
                  t_{2.5\%} \approx 1.97.
                \]
                Dado que \(|3.0| > 1.97\), \textbf{rechazamos} \( H_0 \).
                Existe evidencia suficiente para afirmar que \( \beta_1 \) es
                significativamente distinto de cero al 5\% de significancia.

          \item El intervalo de confianza al 95\% es:
                \[
                  \left[
                    \hat\beta_1 - t_{2.5\%} \times s.e.(\hat\beta_1),
                    \hat\beta_1 + t_{2.5\%} \times s.e.(\hat\beta_1)
                  \right]
                  =
                  \left[
                    2.7 - 1.97 \times 0.9;
                    2.7 + 1.97 \times 0.9
                  \right]
                  \\=            
                  \left[
                    0.927, 4.473
                  \right].
                \]
                
                Con 95\% de confianza, se espera que
                el verdadero efecto de \( X_1 \) sobre la variable dependiente
                esté entre 0.927 y 4.473 unidades. Es decir, un aumento en
                \( X_1 \) está asociado, en promedio, a un aumento {\bf positivo} y
                significativo en la variable dependiente.
        \end{enumerate}}

\item Se tiene un modelo de regresión lineal de la forma:
$$y=\beta_0 + \beta_1 x_1+\beta_2x_2+\beta_3 x_3 + u$$
los resultados de la regresión con 250 datos son los siguientes:
        \[
          \hat{\beta} = \begin{bmatrix} 3 \\ 9 \\ 5 \\2 \end{bmatrix},
        \]
y 
        \[
          (X'X)^{-1} =
          \begin{bmatrix}
1 & 0  & 0  &  0\\
0 & 1  & 0  & -1\\
0 & 0  & 2  & -1\\
0 & -1 & -1 &  1
\end{bmatrix},
          \quad \hat{\sigma}^2 = 2
        \]
\begin{enumerate}
    \item Plantee la hipótesis que testea que $\beta_1=4\beta_3$.

    \item ¿Se puede rechazar la hipótesis nula con la información que disponemos con una significancia del 5\%?

\end{enumerate}

{\color{red}
\textbf{Solución}
\begin{enumerate}
    \item $$H_0 : \beta_1-4\beta_3 = 0$$
$$H_1 : \beta_1-4\beta_3 \neq 0$$
\item el estadistico t esta dado por:
$$t = \frac{\hat{\beta}_1-4\hat{\beta}_3}{s.e.(\hat{\beta}_1-4\hat{\beta}_3)}$$

$$s.e.(\hat{\beta}_1-4\hat{\beta}_3)=\sqrt{\hat{Var}(\hat{\beta}_1-4\hat{\beta}_3)}=\sqrt{\hat{Var}(\hat{\beta}_1)+16\hat{Var}(\hat{\beta}_3)-8\hat{Cov}(\hat{\beta}_1,\hat{\beta}_3)}$$
$$s.e.(\hat{\beta}_1-4\hat{\beta}_3)=\sqrt{2+16\cdot 2-8 \cdot (-2)}=\sqrt{50}$$
por tanto 
$$t = \frac{9-4\cdot2}{7,071}=0.141$$
y con el valor $t_{1-\alpha/2,246}=1.97$. No se rechaza $H_0$.
\end{enumerate}}

\item Considere la siguiente regresión lineal múltiple estimada con \(n = 75\) observaciones:

\[
y \;=\; \beta_{0} + \beta_{1}x_{1} + \beta_{2}x_{2} + \beta_{3}x_{3} + \beta_{4}x_{4} + u.
\]

El modelo arroja un coeficiente de determinación
\[
R^2 = 0.62.
\]

El investigador sospecha que las variables \(x_{3}\) y \(x_{4}\) no aportan capacidad explicativa.  
Para contrastar su hipótesis estima el siguiente modelo:

\[
y \;=\; \gamma_{0} + \gamma_{1}x_{1} + \gamma_{2}x_{2} + e,
\qquad\text{obteniendo}\quad
R^2 = 0.55.
\]
Formule y realice un test de hipótesis que permita decidir, con un 5 \% de significancia, si las variables \(x_{3}\) y \(x_{4}\) son relevantes para el modelo. ($F_{0.95,(2,70)}=3.13$)


{\color{red} \textbf{Solución}

\textbf{Hipótesis:}
\[
H_0:\ \beta_3=\beta_4=0
\quad\text{vs}\quad
H_1:\ \text{al menos una de } \beta_3,\beta_4 \neq 0.
\]
Número de restricciones: $q=2$. En el modelo no restringido hay $k=4$ regresores (sin contar la constante), por lo que
\[
\text{gl}_1=q=2, \qquad \text{gl}_2=n-k-1=75-4-1=70.
\]


\textbf{Estadístico de prueba (basado en $R^2$):}
\[
F=\frac{(R_u^2-R_r^2)/q}{(1-R_u^2)/(n-k-1)}
=\frac{(0.62-0.55)/2}{(1-0.62)/(75-4-1)}
=\frac{0.07/2}{0.38/70}
=6.447.
\]

\medskip
\textbf{Regla de decisión:} Rechazar $H_0$ si $F>F_{0.95,(2,70)}=3.13$.

\medskip
\textbf{Decisión:} Como $6.447>3.13$, se \emph{rechaza} $H_0$ al 5\%.

\medskip
\textbf{p\hbox{-}valor:}
\[
p\text{-valor}=P\!\left(F_{2,70}\ge 6.447\right)\approx 0.0027<0.05.
\]

\medskip
\textbf{Conclusión:} Con evidencia al 5\% (incluso al 1\%), $x_3$ y $x_4$ son conjuntamente \emph{relevantes} para explicar $y$; deben permanecer en el modelo.}
\end{enumerate}

\end{document}